\documentclass{article}
\usepackage[T2A]{fontenc}
\usepackage[utf8]{inputenc}   
\usepackage[english, russian]{babel}

% Set page size and margins
% Replace `letterpaper' with`a4paper' for UK/EU standard size
\usepackage[a4paper,top=2cm,bottom=2cm,left=2cm,right=2cm,marginparwidth=1.75cm]{geometry}

\usepackage{amsmath}
\usepackage{graphicx}
\usepackage[colorlinks=true, allcolors=blue]{hyperref}
\usepackage{amsfonts}
\usepackage{amssymb}
% \usepackage[left=1cm,right=1cm,top=1cm,bottom=1cm]{geometry}
\usepackage{hyperref}
\usepackage{seqsplit}
\usepackage[dvipsnames]{xcolor}
\usepackage{enumitem}
\usepackage{algorithm}
\usepackage{algpseudocode}
\usepackage{algorithmicx}
\usepackage{mathalfa}
\usepackage{mathrsfs}
\usepackage{dsfont}
\usepackage{caption,subcaption}
\usepackage{wrapfig}
\usepackage[stable]{footmisc}
\usepackage{indentfirst}
\usepackage{rotating}
\usepackage{pdflscape}
\usepackage{listings}
\usepackage{amsmath}
\usepackage{graphicx}
\usepackage[colorlinks=true, allcolors=blue]{hyperref}
\usepackage{amsfonts}
\usepackage{amssymb}
% \usepackage[left=1cm,right=1cm,top=1cm,bottom=1cm]{geometry}
\usepackage{hyperref}
\usepackage{seqsplit}
\usepackage[dvipsnames]{xcolor}
\usepackage{enumitem}
\usepackage{algorithm}
\usepackage{algpseudocode}
\usepackage{algorithmicx}
\usepackage{mathalfa}
\usepackage{mathrsfs}
\usepackage{dsfont}
\usepackage{caption,subcaption}
\usepackage{wrapfig}
\usepackage[stable]{footmisc}
\usepackage{indentfirst}
\usepackage{rotating}
\usepackage{pdflscape}
\usepackage{graphicx}

\usepackage{minted}

\usepackage{MnSymbol,wasysym}
\usepackage{tikz}

\usepackage{minted}

\usepackage{MnSymbol,wasysym}

\title{Домашнее задание 10}
\author{Лиза Фоменко}
\date{}

\begin{document}

\maketitle

\textbf{Турнир} - это полный ориентированный граф, то есть такой ориентированный граф, в котором между любой парой различных вершин есть ровно одно ребро.

\section*{\textbf{Задание 1}} \textbf{\textit{Докажите, что в турнире на $n$ вершинах есть простой (несамопересекающийся) путь длины $n - 1$. Постройте алгоритм, находящий такой путь, и оцените время его работы.}}

Попробуем доказать индуктивно.

\textbf{База индукции:} для n = 2 очевидно: у нас есть ребро $v_1$ -> $v_2$, то есть путь длины $n-1 = 1$. 

\textbf{Предположение индукции:} Пусть для турнира с $n-1$ вершиной верно, что существует простой путь длины $n-2$ 

\textbf{Шаг индукции:} Рассмотрим турнир из $n$ вершин. Заметим, что если убрать из него случайную вершину $v$ вместе с выходящими/входящими в нее ребрами, то мы получим турнир из $n-1$ вершин, для которого известно, что существует такой путь из $n-2$ вершин вида $v_1, ... v_{n-1}$. 

Попробуем вернуть вершину $v$ в граф. Пойдем от вершины $v_1$ по имеющемуся простому циклу длины $n-2$, то есть $v_1, v_2, v_3, ...$. Для каждой вершины проверяем, существует ли ребро $v, v_i$ ("вход" $v$ в цикл): заметим, что если его не существует, то обязательно существует $v_i, v$ ("выход" из цикла в $v$). Найдем первую вершину, для которой существует  $v, v_i$ - через него "$v$ войдет в цикл". Так как мы проверяем каждую вершину по циклу, то для всех вершин, находящихся раньше найденной $v_i$, существует ребро, ведущее в $v$. Тогда мы можем "убрать" из нашего пути ребро $v_{i-1}, v_i$ и добавить в него два ребра: $v_{i-1}, v$ и $v, v_i$. Тогда вершина $v$ войдет в путь, а его длина увеличится на 1: получим путь из $n$ вершин длины $n-1$ - простой путь.

Перейдем к алгоритму поиска. На примере доказательтсва по индукции мы можем точно так же построить алгоритм поиска, добавляя по одной вершине за ход алгоритма.

\begin{verbatim}

G = (V, E)            # наш граф. здесь не оговорено как он задан


v0 = G[0]             # пример выбора начальной вершины
v1 = v0_merged[0]     # выбираем вторую вершину: одна из тех, в которые существует ребро из v0: v0->v1

path = [v0, v1]

for v in V - v0 - v1:                                # для всех оставшихся вершин
    if v not in path:                                # если она еще не входит в путь
        for i in range(len(path)):                   # идем по циклу
            if v -> path[i]:                         # если существует путь v -> v_i
                path = path[:i] + v + path[i:]       # включаем ее в цикл по правилу выше
                break
                
return path                
\end{verbatim}

Этот алгоритм работает за $O(n^2):$ мы проходим цикл для каждой вершины ($n-2, O(n)$), и просматриваем вершины из уже существующего пути, что тоже занимает $O(n) \Rightarrow $ асимптотика $O(n^2)$. Его корректность доказана при доказательстве существования самого пути :) 

\bigskip

\section*{\textbf{Задание 2}} \textbf{\textit{Примените DFS к графу, начиная с вершины $e$. Найдите время открытия и закрытия каждой вершины, приведите лес обхода в глубину. Для рисования графов рекомендуется использовать \href{https://madebyevan.com/fsm/}{вот эту штуку}}}.


Выпишем порядок вершин: a, b, c, d, e, f. Обход начинается в $e$, идем в $d$, из нее никуда пойти не можем => возвращаемся в $e$ и проходим в $f$, из которой путей тоже не существует. Получили первое дерево обхода.

Для второго дерева обхода начнем с вершины $a$ (так как идем в алфавитном порядке), затем идем через $b$ в $c$. Второе дерево обхода готово. 

У меня люто висел сайт который был рекомендован, поэтому я вместо борьбы с латехом нарисовала картинки за 5 минут и прикрепляю их в виде .png сюда)

\begin{figure}[H]
     \centering
     \begin{subfigure}[b]{0.45\textwidth}
         \centering
         \includegraphics[width=\textwidth]{algos/dfs_times.png}
         \caption{DFS times}
     \end{subfigure}
     \hfill
     \begin{subfigure}[b]{0.45\textwidth}
         \centering
         \includegraphics[width=\textwidth]{algos/dfs_tree.png}
         \caption{DFS tree}
     \end{subfigure}
     \hfill
\end{figure}

\bigskip


\section*{\textbf{Задание 3}} \textbf{\textit{Постройте алгоритм, находящий все стоки в ориентированном графе, если известно, что их $\Theta(\sqrt{|V|})$. Сток - это вершина с нулевой исходящей степенью. Граф уже лежит в памяти в виде матрицы смежности.}}

Пусть в матрице исходящие вершины находятся в строке (чтобы оперировать можно было количеством столбцов). 
Наивный алгоритм работает за $O(n^2)$ - при просмотре каждого элемента матрицы смежности.
Но тут никак не учитывается количество стоков. Поэтому попробуем ускорить алгоритм, используя это великое знание.

Идея: делим матрицу на две части по количеству столбцов: $\alpha$ и $n-\alpha$. В первой итерации алгоритма мы будем рассматривать только первые $\alpha n$ столбцов. Также нам нет никакой необходимости идти по всем столбца для вершины, если мы уже встретили 1 в матрице в этом ряду. 
Собственно что мы будем делать: для всех $|V|=n$ вершин смотрим первые $\alpha n$ столбцов. Идем по столбцам для каждой вершины, если встречаем 1, выкидываем вершину из кандидатов в стоки. 
После полного прохода, который займет $O(n * \alpha)$, у нас останется набор вершин-кандидатов в стоки размера $n'$, для которых на первых $\alpha$ столбцах находятся только 0. 
Далее только для этих вершин мы будем аналогичным образом рассматривать $n - \alpha$ столбцов. Этот этап у нас займет $O(n' * (n-\alpha))$. 
Итоговая асимптотика: $O(n * \alpha) + n' * (n-\alpha)$
Мы знаем, что она должна быть меньше $O(n^2)$, при этом $n$ в явном виде входит в оба слагаемых. Тут я приведу рассуждения, которые не могу обосновать математически, но они мне кажутся по наитию верными.
У нас после первого шага остается $n'$ кандидатов. Мы знаем, что всего стоков $\Theta(\sqrt{|V|}) = \Theta(\sqrt{n})$. Тогда мы можем грубо оценить $n'$ через $O(\sqrt{n})$, и тогда асимптотика второй части алгоритма будет $O(\sqrt{n} * (n-\alpha))$, что приводит нас к $O(n\sqrt{n})$. Можно ли подобрать $\alpha$ так, чтобы первая часть алгоритма тоже была $O(n\sqrt{n})$? Можно, при $\alpha = \sqrt{n}$. Тогда итоговая сложность составляет $O(n\sqrt{n})$.

Можно ли быстрее? Могут ли обе части работать быстрее? Мне кажется нет, потому что во второй части алгоритма у нас всегда будет зависимость от $\Theta(\sqrt{n})$. При делении столбцов у нас всегда останется $1 - \alpha$ столбцов, что всегда при $\alpha << n$ будет линией. При $\alpha$ ~ $n$ мы получим $O(n^2)$ в первой части алгоритма. 

\bigskip

\end{document}